
% =============================================================================
% Project : THE_PHILOSOPHY_OF_VIABILITY
% Book    : Book I — Ontology of Viability
% Book ID : VS.B01
% Part    : Part I — Identity and Manifestation
% Part ID : VS.B01.P01
% Chapter : Chapter 08 — Comparative Philosophical Examination
% Chapter ID : VS.B01.P01.C08
%
% File :
%   4_ASSETS/diagrams/
%   VS.B01.P01.C08.S13.SS08.P005.DG01.avicenna_modal_causal_dependence.tex
%
% Asset :
%   DG01 — Avicenna's Modal and Causal Dependence
%
% Structural Identifier :
%   VS.B01.P01.C08.S13.SS08.P005.DG01
%
% Purpose :
%   To distinguish the possible existent considered in itself
%   from the same existent considered as necessary through
%   its complete cause, and to indicate the explanatory
%   dependence of caused existents upon the First Principle.
%
% Dependencies :
%   tikz, positioning, arrows.meta.
%   These are already loaded by the Book I preamble.
%
% Philosophical and Representational Limits :
%   1. The diagram is schematic and not an exhaustive map
%      of Avicenna's cosmology.
%   2. The two modal descriptions concern the same caused
%      existent under different respects; they are not
%      successive temporal stages.
%   3. The arrows indicate causal or explanatory dependence,
%      not temporal succession.
%   4. The causal series is simplified for explanatory clarity.
%   5. The diagram does not imply that the First Principle
%      is identical with Being as such or First Determination.
%   6. The diagram does not map Avicenna's cosmology onto
%      the proposed realms and hubs of Book I.
%   7. Causal dependence is not equated with Manifestation.
%
% =============================================================================

\begin{figure}[htbp]
    \centering

    \begin{tikzpicture}[
        node distance=0.78cm,
        principle/.style={
            draw=ViabilityGold,
            fill=ViabilityPaleGold,
            rounded corners=2pt,
            line width=0.8pt,
            text width=0.70\linewidth,
            align=center,
            inner sep=7pt
        },
        mediation/.style={
            draw=ViabilityBlue,
            fill=ViabilityPaleBlue,
            rounded corners=2pt,
            line width=0.7pt,
            text width=0.70\linewidth,
            align=center,
            inner sep=7pt
        },
        caused/.style={
            draw=ViabilityGreen,
            fill=ViabilityPaleGreen,
            rounded corners=2pt,
            line width=0.8pt,
            text width=0.82\linewidth,
            align=center,
            inner sep=8pt
        },
        dependence/.style={
            -{Stealth[length=2.2mm]},
            draw=ViabilityDarkBlue,
            line width=0.85pt
        },
        relationlabel/.style={
            fill=white,
            text=ViabilityDarkBlue,
            font=\footnotesize,
            inner sep=2pt,
            align=center
        }
    ]

    \node[principle] (first) {
        \textbf{Necessary Existent in Itself}\\[3pt]
        Uncaused; independent; not necessitated by another
    };

    \node[mediation, below=of first] (series) {
        \textbf{Causal Order}\\[3pt]
        Intermediate causes and relations of dependence
    };

    \node[mediation, below=of series] (cause) {
        \textbf{Complete Cause of a Caused Existent}\\[3pt]
        The causal conditions sufficient for its existence
    };

    \node[caused, below=of cause] (effect) {
        \textbf{One Caused Existent Considered under Two Respects}\\[7pt]

        \begin{tabular}{c}
            \textbf{In itself:}\\
            Possible Existent
            \\[5pt]
            \textbf{In relation to its complete cause:}\\
            Necessary through Another
        \end{tabular}
    };

    \draw[dependence]
        (first) --
        node[relationlabel, right] {ultimate causal\\grounding}
        (series);

    \draw[dependence]
        (series) --
        node[relationlabel, right] {mediated causal\\dependence}
        (cause);

    \draw[dependence]
        (cause) --
        node[relationlabel, right] {necessitation\\through another}
        (effect);

    \end{tikzpicture}

    \caption[
        Avicenna's distinction between possibility in itself
        and necessity through another
    ]{
        Avicenna's modal and causal account of a caused existent.
        The possible existent remains possible when considered
        in itself, while its existence is necessary through
        another when considered in relation to its complete
        cause. The causal order is represented schematically;
        the diagram does not depict a temporal sequence.
    }

    \label{fig:VS.B01.P01.C08.S13.SS08.P005.DG01}

\end{figure}

% =============================================================================
% END OF DIAGRAM DG01
% =============================================================================
